\chapter{Adenoidectomy}

\section*{Introduction \& Clinical Indications}
Adenoidectomy is a common pediatric surgical procedure involving the complete or partial removal of the adenoids, which are a mass of lymphoid tissue located on the posterior superior wall of the nasopharynx, behind the soft palate. This surgery is primarily indicated when hypertrophied (enlarged) adenoids cause significant obstruction of the nasal airway, leading to chronic mouth breathing, snoring, and sleep-disordered breathing, or when they contribute to recurrent infections such as otitis media, adenoiditis, and sinusitis. By removing this obstructing and often infected tissue, the procedure aims to restore normal nasal breathing, improve sleep quality, and reduce the frequency of upper respiratory and middle ear pathologies, thereby enhancing a child's overall health and development.

\begin{itemize}
  \item Adenoid removal
  \item Pharyngeal tonsillectomy
  \item Excision of adenoid vegetations
  \item Nasopharyngeal lymphoid tissue ablation
  \item Adenoidectomy with tonsillectomy (when performed concurrently)
  \item Isolated adenoidectomy (when performed alone)
\end{itemize}

The fundamental principle of adenoidectomy is the mechanical removal of hypertrophied lymphoid tissue that is causing symptomatic obstruction or acting as a chronic reservoir for infection. The adenoids, while part of the immune system, can become pathologically enlarged, leading to a cascade of issues. Obstruction of the nasopharynx directly impairs nasal breathing, forcing mouth breathing, which can lead to dental malocclusion, altered facial growth, and sleep-disordered breathing, including obstructive sleep apnea. Furthermore, the proximity of the adenoids to the Eustachian tube openings means that enlargement can block these tubes, preventing proper middle ear ventilation and drainage, thereby predisposing to chronic otitis media with effusion and recurrent acute otitis media.

The rationale for surgical intervention is to alleviate these symptoms by physically excising the offending tissue. This restores patency to the nasal airway, allowing for normal nasal respiration and improving sleep quality. It also re-establishes normal Eustachian tube function, reducing the incidence of middle ear disease. Technically, the goals are complete or near-complete removal of the adenoid mass while meticulously preserving the surrounding vital structures, particularly the Eustachian tube orifices and the soft palate, and achieving meticulous hemostasis to prevent postoperative bleeding. By removing the chronically inflamed or infected tissue, the procedure also aims to reduce the bacterial load and inflammatory response in the nasopharynx.

The adenoids were first systematically described by the Danish physician Wilhelm Meyer in 1868. Meyer meticulously documented the clinical syndrome associated with adenoid hypertrophy, including nasal obstruction, hearing loss, and characteristic facial features, and subsequently pioneered their surgical removal, which he termed the treatment for "adenoid vegetations." His work marked a significant milestone, establishing adenoidectomy as a recognized and effective surgical intervention.

Throughout the late 19th and early 20th centuries, adenoidectomy, often combined with tonsillectomy, gained widespread acceptance. Early techniques primarily involved blind curettage using instruments like the Beckmann adenoid curette, which, while effective, carried risks of incomplete removal or injury to adjacent structures. The mid-20th century saw refinements in surgical technique, including improved visualization with mirrors and later endoscopes, leading to more precise tissue removal and better hemostasis. The advent of general anesthesia and improved surgical instruments further enhanced safety and efficacy. Modern techniques, such as powered microdebriders, radiofrequency ablation (coblation), and electrocautery, allow for highly controlled tissue removal under direct endoscopic visualization, minimizing trauma and improving outcomes. The indications for surgery have also evolved, moving from a more liberal approach to evidence-based criteria focusing on significant obstruction, sleep-disordered breathing, and recurrent infections.

Adenoidectomy is indicated for a range of conditions primarily related to adenoid hypertrophy and chronic inflammation. The decision for surgery is typically based on persistent symptoms despite conservative management.

\begin{itemize}
  \item **Chronic nasal obstruction:** Persistent mouth breathing, snoring, and nasal congestion that significantly impair quality of life and are unresponsive to medical therapy.
  \item **Obstructive sleep apnea (OSA) in children:** Documented sleep-disordered breathing, often confirmed by polysomnography, where adenoid hypertrophy is a primary contributing factor.
  \item **Recurrent acute otitis media (RAOM):** Frequent episodes of middle ear infection (typically 3 or more in 6 months, or 4 or more in 12 months) that are refractory to antibiotic treatment, especially when associated with adenoiditis.
  \item **Chronic otitis media with effusion (COME):** Persistent fluid in the middle ear for 3 months or longer, leading to hearing loss, particularly when associated with adenoid hypertrophy and often combined with tympanostomy tube insertion.
  \item **Chronic adenoiditis:** Persistent inflammation and infection of the adenoids, leading to chronic rhinorrhea, post-nasal drip, and halitosis.
  \item **Recurrent or chronic sinusitis:** Adenoid hypertrophy can impede sinus drainage and act as a bacterial reservoir, contributing to recurrent or chronic sinus infections.
  \item **Speech abnormalities:** Rarely, severe adenoid hypertrophy can contribute to hyponasal speech (muffled voice) or other speech impediments.
  \item **Orofacial development issues:** Long-standing mouth breathing due to adenoid hypertrophy can influence dental occlusion and facial growth patterns.
\end{itemize}

Adenoidectomy is considered a primary surgical treatment for significant adenoid hypertrophy and its associated complications. It is often the definitive intervention when conservative medical management, such as allergy treatment, nasal steroids, or antibiotics for acute infections, has failed to alleviate symptoms. For children with moderate to severe obstructive sleep apnea, adenoidectomy (often combined with tonsillectomy) is typically the first-line surgical approach, demonstrating high success rates in resolving or significantly improving symptoms.

While it is a primary treatment, it can also be viewed as a secondary option in the sense that a period of observation or medical therapy is usually attempted first for milder symptoms. For instance, in cases of chronic otitis media with effusion, adenoidectomy is often considered after a period of watchful waiting or after initial tympanostomy tube insertion if effusions recur. Thus, its positioning is dynamic: it is the primary surgical solution for established, significant pathology, but it is typically pursued after less invasive measures have been explored or proven insufficient.

\section*{Relevant Surgical Anatomy}
The adenoids, also known as the pharyngeal tonsil, are a collection of lymphoid tissue situated in the nasopharynx, specifically on its posterior superior wall, extending from the sphenoid bone to the level of the soft palate. They are part of Waldeyer's ring, a circular arrangement of lymphoid tissue that guards the entrance to the pharynx, comprising the palatine tonsils, lingual tonsils, and tubal tonsils. The adenoids are typically largest between the ages of 3 and 7 years and usually atrophy after puberty.

Key anatomical relations include the posterior choanae (the openings from the nasal cavity into the nasopharynx) inferiorly, and the openings of the Eustachian tubes (connecting the middle ear to the nasopharynx) laterally. Hypertrophy of the adenoids can directly obstruct the choanae, impeding nasal airflow, and can also impinge upon the Eustachian tube orifices, leading to middle ear ventilation dysfunction and recurrent otitis media. The blood supply to the adenoids is rich, primarily derived from branches of the ascending pharyngeal artery, ascending palatine artery (from the facial artery), sphenopalatine artery (from the maxillary artery), and the pharyngeal branch of the internal maxillary artery. Venous drainage occurs via the pharyngeal plexus, which empties into the internal jugular vein. Lymphatic drainage is to the retropharyngeal and deep cervical lymph nodes. The nerve supply is from the pharyngeal plexus, formed by branches of the glossopharyngeal and vagus nerves. Surgical landmarks include the posterior edge of the hard palate and the Eustachian tube orifices, which must be carefully avoided during tissue removal.

\section*{Preoperative Workup \& Preparation}
A comprehensive preoperative workup is essential to ensure patient safety and optimize surgical outcomes. This typically begins with a detailed history and physical examination, focusing on the child's general health, cardiorespiratory status, and any history of bleeding disorders in the patient or family. Specific attention is paid to the severity and duration of symptoms related to adenoid hypertrophy.

\begin{itemize}
  \item **Medical Evaluation:** A thorough medical history is taken, including allergies, current medications, and previous anesthetic experiences. A physical examination includes an assessment of the airway, heart, and lungs. For children with significant comorbidities, a pediatric consultation or cardiology clearance may be required.
  \item **Laboratory Tests:** Routine preoperative blood tests may include a complete blood count (CBC) to check for anemia or infection, and a coagulation profile (PT/INR, PTT) to screen for bleeding disorders, especially if there is a suspicious history.
  \item **Imaging/Diagnostic Studies:** While not always routine, a lateral neck X-ray can confirm adenoid hypertrophy and assess the nasopharyngeal airway space. Flexible nasopharyngoscopy can provide direct visualization of the adenoids and their relationship to the Eustachian tubes and choanae. Polysomnography is indicated for children with suspected obstructive sleep apnea.
  \item **Medication Adjustment:** Medications that increase bleeding risk, such as aspirin, NSAIDs (ibuprofen, naproxen), and herbal supplements, should be discontinued 1-2 weeks prior to surgery.
  \item **NPO Status:** Patients must adhere to strict NPO (nil per os) guidelines before surgery to prevent aspiration during anesthesia, typically no solid food for 6-8 hours and clear liquids for 2-4 hours.
  \item **Prophylactic Antibiotics:** Prophylactic antibiotics are generally not routinely administered for isolated adenoidectomy but may be considered in specific high-risk cases or when combined with other procedures.
  \item **Infection Control:** Any acute upper respiratory tract infection or active adenoiditis should be treated and resolved before elective surgery to minimize anesthetic and surgical risks.
  \item **Informed Consent:** Detailed discussion with parents/guardians regarding the procedure, its benefits, risks, alternatives, and expected postoperative course is crucial, obtaining fully informed consent.
\end{itemize}

Careful patient selection is paramount for adenoidectomy. Certain conditions represent absolute contraindications, while others require significant precautions.

\textbf{Absolute Contraindications:}
\begin{itemize}
  \item **Unrepaired cleft palate or known submucous cleft palate:** Adenoidectomy in these patients can worsen velopharyngeal insufficiency (VPI), leading to hypernasal speech and nasal regurgitation of food and liquids, as the adenoids contribute to velopharyngeal closure.
  \item **Known or suspected bleeding disorder (uncorrected):** Conditions like hemophilia, von Willebrand disease, or severe thrombocytopenia pose an unacceptable risk of life-threatening hemorrhage unless fully diagnosed and managed preoperatively.
  \item **Acute infection of the pharynx or nasopharynx:** Active bacterial or viral infections can increase the risk of anesthetic complications, postoperative infection, and bleeding. Surgery should be deferred until the infection has resolved.
\end{itemize}

\textbf{Relative Contraindications \& Precautions:}
\begin{itemize}
  \item **Velocardiofacial syndrome (22q11.2 deletion syndrome):** These patients often have subtle palatal abnormalities and are at higher risk for VPI post-adenoidectomy. Careful preoperative evaluation and discussion with a craniofacial team are necessary.
  \item **Neuromuscular disorders:** Children with conditions affecting swallowing or palatal function (e.g., cerebral palsy) may be at increased risk for VPI or aspiration post-surgery.
  \item **Very young age (e.g., under 1 year):** While not an absolute contraindication, adenoidectomy in infants is less common and requires careful consideration due to smaller airway size and potential for other underlying causes of obstruction.
  \item **Severe systemic illness or uncontrolled medical conditions:** Patients with unstable cardiac, pulmonary, or metabolic diseases may have increased anesthetic and surgical risks, necessitating optimization of their medical status before surgery.
  \item **Anatomical abnormalities of the cervical spine:** Rarely, conditions like atlantoaxial instability (e.g., in Down syndrome) require careful positioning to avoid spinal cord injury during intubation and surgery.
\end{itemize}

\section*{Surgical Technique \& Procedure}
Adenoidectomy is performed under general anesthesia, typically as an outpatient procedure. The patient is positioned supine with the neck extended, often using a shoulder roll, to optimize exposure of the nasopharynx. A mouth gag (e.g., Boyle-Davis or Crowe-Davis) is inserted to keep the mouth open and depress the tongue.

The key operative steps are as follows:
\begin{itemize}
  \item **Exposure:** The soft palate is gently retracted superiorly, often using a red rubber catheter passed through the nose and out the mouth, or a specialized palatal retractor, to provide a clear view of the nasopharynx and the adenoid pad. Endoscopic visualization via a rigid or flexible endoscope passed transorally or transnasally is increasingly common to ensure complete removal and minimize injury.
  \item **Adenoid Removal:** The adenoid tissue can be removed using several techniques:
    \begin{itemize}
      \item **Curettage:** A specialized adenoid curette (e.g., Beckmann) is passed behind the soft palate and used to scrape the adenoid tissue from the posterior nasopharyngeal wall. This is often performed blindly or with mirror guidance.
      \item **Powered Microdebrider:** A rotating blade within a protective sheath is used to precisely shave and suction the adenoid tissue under direct endoscopic visualization. This method allows for controlled removal and preservation of surrounding structures.
      \item **Radiofrequency Ablation (Coblation):** Uses a low-temperature radiofrequency energy to dissolve tissue, offering precise removal with minimal thermal injury and good hemostasis.
      \item **Electrocautery:** Monopolar or bipolar electrocautery can be used to excise and coagulate the adenoid tissue, providing excellent hemostasis.
    \end{itemize}
  \item **Hemostasis:** Meticulous control of bleeding from the adenoid bed is crucial. This is achieved using suction cautery, bipolar cautery, or by applying pressure with a gauze pack soaked in vasoconstrictive agents (e.g., epinephrine) for several minutes.
  \item **Inspection:** The nasopharynx is thoroughly inspected, often endoscopically, to ensure complete removal of hypertrophied tissue, confirm hemostasis, and verify that the Eustachian tube orifices and posterior choanae are clear and undamaged.
  \item **Concomitant Procedures:** If indicated, a tonsillectomy or tympanostomy tube insertion is performed during the same anesthetic.
  \item **Closure:** No formal closure is required for the adenoid bed. The mouth gag is removed, and the patient is awakened from anesthesia.
\end{itemize}
The procedure typically lasts 15-30 minutes for an isolated adenoidectomy. Most children are discharged home after a short period of observation in the post-anesthesia care unit (PACU).

\section*{Complications \& Risks}
While adenoidectomy is generally a safe procedure, like all surgeries, it carries potential risks and complications. These can be broadly categorized into general surgical risks and procedure-specific risks.

\textbf{General Surgical Complications:}
\begin{itemize}
  \item **Anesthesia risks:** Reactions to anesthetic agents, respiratory depression, aspiration, and cardiac events.
  \item **Bleeding:** Postoperative hemorrhage, though usually minor, can occasionally be significant enough to require re-operation or blood transfusion.
  \item **Infection:** Wound infection, though uncommon in the nasopharynx, or systemic infection.
  \item **Pain:** Postoperative discomfort, usually mild and manageable.
  \item **Nausea and vomiting:** Common after general anesthesia.
\end{itemize}

\textbf{Procedure-Specific Complications:}
\begin{itemize}
  \item **Velopharyngeal insufficiency (VPI):** The most significant specific risk, particularly in children with undiagnosed submucous cleft palate or other palatal abnormalities. It results in hypernasal speech and nasal regurgitation of fluids, which may be temporary or, rarely, permanent, requiring speech therapy or further surgical correction.
  \item **Injury to Eustachian tube orifices:** Can lead to persistent middle ear dysfunction or scarring, though rare with modern techniques.
  \item **Nasopharyngeal stenosis:** Extremely rare, but excessive scarring can lead to narrowing of the nasopharynx.
  \item **Changes in voice quality:** Temporary or, rarely, permanent changes in voice, often a more "open" or less muffled sound, which is usually perceived as an improvement but can occasionally be hypernasal.
  \item **Regrowth of adenoid tissue:** In a small percentage of cases, adenoid tissue can regrow, leading to recurrence of symptoms and potentially requiring revision surgery.
  \item **Dental injury:** Damage to teeth from the mouth gag, especially in children with loose primary teeth or dental anomalies.
  \item **Atlantoaxial subluxation:** An extremely rare but severe complication, particularly in patients with underlying conditions like Down syndrome, due to hyperextension of the neck during surgery.
  \item **Incomplete relief of symptoms:** Despite successful surgery, some children may not experience complete resolution of symptoms, especially if other factors contribute to their condition (e.g., allergic rhinitis, obesity in OSA).
\end{itemize}

\section*{Postoperative Care \& Recovery}
The recovery from adenoidectomy is generally rapid and well-tolerated, especially when performed as an isolated procedure.

Immediately after surgery, the child is monitored in the Post-Anesthesia Care Unit (PACU) for vital signs, airway patency, and signs of bleeding. Once awake, alert, and stable, with pain controlled and able to tolerate oral fluids, most children are discharged home within a few hours. The first 24-48 hours typically involve mild throat and nasal discomfort, which is usually managed with over-the-counter pain relievers like acetaminophen or ibuprofen. A low-grade fever is common. Parents are encouraged to offer plenty of clear fluids to maintain hydration and prevent dehydration. A soft, bland diet is recommended for the first day or two, gradually progressing to a regular diet as tolerated.

Over the first 1-2 weeks, the child should avoid strenuous activities, heavy lifting, and contact sports to minimize the risk of bleeding. School attendance can usually resume within 3-7 days, depending on the child's comfort level. Nasal congestion and a mild nasal discharge are common as the surgical site heals. Any changes in voice quality, such as a more open or slightly nasal sound, usually resolve within a few weeks as the child adapts. Long-term recovery involves the gradual resolution of chronic symptoms like mouth breathing, snoring, and recurrent infections, with the full benefits often becoming apparent over several weeks to months.

During adenoidectomy, the surgeon typically documents the size, consistency, and degree of obstruction caused by the adenoid tissue. Findings may include diffuse hypertrophy, often described as grade I-IV based on the percentage of choanal obstruction, or signs of chronic inflammation such as erythema, edema, and mucopurulent discharge. The presence of any unusual asymmetry, masses, or friable tissue is noted.

The resected adenoid tissue is routinely sent for histopathological examination. In the vast majority of cases, the pathology report confirms benign lymphoid hyperplasia, characterized by an increase in the number and size of lymphoid follicles with active germinal centers, consistent with a normal immune response to chronic antigenic stimulation. This finding confirms that the enlargement was reactive rather than neoplastic. However, in rare instances, particularly if there were atypical preoperative findings or unusual macroscopic appearance, the pathologist will specifically look for other conditions such as lymphoma, nasopharyngeal carcinoma, or juvenile nasopharyngeal angiofibroma. The pathological confirmation of benign lymphoid hyperplasia provides reassurance and correlates with the clinical picture of chronic inflammation and hypertrophy.

A normal and successful postoperative course following adenoidectomy is characterized by a significant improvement in the child's symptoms. Pain is typically mild and resolves within a few days, often managed effectively with non-opioid analgesics. Nasal breathing should improve almost immediately, leading to a reduction or cessation of mouth breathing and snoring. Sleep quality is expected to improve, with a decrease in episodes of obstructive sleep apnea or sleep-disordered breathing. Any associated nasal discharge or congestion from chronic adenoiditis should gradually resolve.

Over the subsequent weeks and months, parents can expect a noticeable reduction in the frequency and severity of recurrent acute otitis media and chronic otitis media with effusion, as Eustachian tube function normalizes. The child's overall energy levels, daytime alertness, and potentially even academic performance may improve due to better sleep. The surgical site in the nasopharynx heals completely within a few weeks, with minimal scarring. The long-term expectation is a healthier, more comfortable child with fewer upper respiratory and ear-related health issues.

Postoperative care for adenoidectomy focuses on pain management, hydration, and monitoring for complications. Parents are advised to administer prescribed pain medication as needed, encourage fluid intake, and offer a soft diet initially. Activity restrictions, such as avoiding strenuous exercise, are typically in place for about one week.

\begin{itemize}
  \item **Follow-up Visits:** A routine follow-up appointment with the surgeon is usually scheduled 2-4 weeks after surgery to assess healing, confirm symptom resolution, and address any lingering concerns.
  \item **Suture/Staple Removal:** No sutures or staples are used in the adenoid bed, so no removal is necessary.
  \item **Rehabilitation/Physical Therapy:** Generally not required for adenoidectomy alone. However, if speech abnormalities persist or velopharyngeal insufficiency is suspected, referral to a speech-language pathologist may be indicated.
  \item **Warning Signs:** Parents are educated on critical warning signs that necessitate immediate medical attention. These include significant postoperative bleeding (bright red blood, frequent swallowing), persistent high fever, severe pain unresponsive to medication, difficulty breathing, or refusal to drink fluids leading to signs of dehydration.
\item **Long-term Monitoring:** For children with a history of recurrent otitis media, ongoing audiological assessment and follow-up with an ENT specialist may be necessary to monitor hearing and middle ear status.
\end{itemize}

\section*{Outcomes \& Prognosis}
Adenoidectomy is one of the most frequently performed surgical procedures in pediatric otolaryngology worldwide. The incidence of adenoid hypertrophy and its associated complications peaks in early childhood, typically between the ages of 1 and 8 years, coinciding with the natural growth phase of lymphoid tissue and increased exposure to childhood infections. While both sexes are affected, some studies suggest a slight male predominance for obstructive symptoms. The leading indications for adenoidectomy are obstructive sleep apnea, recurrent acute otitis media, and chronic otitis media with effusion, often necessitating concurrent tonsillectomy or tympanostomy tube insertion. The procedure's morbidity is generally low, with serious complications being rare. Mortality is exceedingly rare, primarily associated with severe anesthetic complications or uncontrolled postoperative hemorrhage. The high prevalence of these childhood conditions contributes to the substantial number of adenoidectomies performed annually, making it a cornerstone of pediatric ENT practice.

The outcomes of adenoidectomy are overwhelmingly positive, with high success rates in resolving the primary indications. For obstructive sleep apnea, adenoidectomy (often with tonsillectomy) leads to significant improvement or resolution of symptoms in 80-90\% of children, as evidenced by studies like the Childhood Adenotonsillectomy Trial (CHAT). In cases of recurrent acute otitis media and chronic otitis media with effusion, adenoidectomy, particularly when combined with tympanostomy tube insertion, significantly reduces the frequency of infections and improves hearing outcomes.

The prognosis is generally excellent, with most children experiencing sustained relief from nasal obstruction, improved sleep quality, and a reduction in ear infections. Recurrence of adenoid hypertrophy requiring revision surgery is uncommon, estimated to occur in less than 5\% of cases, and is more likely in younger children or those with underlying allergic conditions. Prognostic factors for successful outcomes include the absence of underlying craniofacial anomalies, the severity of preoperative symptoms, and the completeness of adenoid tissue removal. While adenoidectomy is highly effective, it is important to note that it does not cure all causes of nasal obstruction or ear problems, and other contributing factors (e.g., allergies, sinusitis) may require ongoing management.

\section*{References}
\begin{enumerate}
  \item Cummings CW, Flint PW, Haughey BH, et al. Cummings Otolaryngology - Head and Neck Surgery, 7th edition. Elsevier; 2021.
  \item Kerschner JE, Preciado D. Adenoidectomy and Tonsillectomy. In: Kliegman RM, St. Geme JW, Blum NJ, et al., editors. Nelson Textbook of Pediatrics, 21st edition. Elsevier; 2020.
  \item Marcus CL, Moore RH, Rosen CL, et al. A randomized trial of adenotonsillectomy for childhood obstructive sleep apnea. New England Journal of Medicine. 2013;368(25):2366-2376.
  \item American Academy of Otolaryngology - Head and Neck Surgery. Clinical Practice Guideline: Otitis Media with Effusion (Update). Otolaryngology--Head and Neck Surgery. 2016;154(1 Suppl):S1-S41.
\end{enumerate}

\clearpage
